% Part: first-order-logic
% Chapter: models-theories
% Section: set-theory

\documentclass[../../../include/open-logic-section]{subfiles}

\begin{document}

\olfileid{fol}{mat}{set}

\olsection{د سټونو تيوري}

د رياضۍ کابو ټولې برخې د سټونو په تيورۍ کښې جوړېدے شي۔ په دې
تيورۍ کښې د رياضۍ جوړول څو کارونه غواړي۔ لومړے، د~$\in$ اړيکې دپاره
د بديهي اصولو يو سټ غواړي۔ بېلابېل بديهي نظامونه جوړ شوي دي چې کله
کله د~$\in$ په باب يو بل سره ټکر لرونکي خاصيتونه لري۔ د~$\Log{ZFC}$
په نوم بديهي نظام، يعنې د انتخاب له بديهي اصل سره د
زرملو--فرېنکل سټ تيوري، ځانګړے مقام لري: دا تر بل هر بديهي نظام ډېر
کارول کېږي او څېړل کېږي، ځکه بديهي اصول ئې د
کابو ټولو هغو خبرو د ثبوت دپاره بس کېږي چې رياضيپوهان ئې ثابتول
غواړي۔ خو د دې له څرګندولو وړاندې بايد لومړے روښانه کړو چې د
رياضيپوهانو مطلوبې خبرې اصلًا څنګه \emph{بيانولے} شو۔ د پيل دپاره،
ژبه هېڅ !!{constant}s يا !!{function}s نۀ لري؛ نو په لومړي نظر دا
روښانه نۀ ده چې د ځانګړو سټونو (لکه~$\emptyset$ يا~$\Nat$)، پر
سټونو د عملونو (لکه~$X \cup Y$ او~$\Pow{X}$)، يا لا د هغو جوړښتونو
په باب څنګه خبرې وکړو چې له سټونو پرته نور څيزونه هم رانغاړي، لکه
اړيکې او تابعې۔

په پيل کښې، ``عنصر دے'' يوازېنۍ مهمه اړيکه نۀ ده؛ ``فرعي سټ دے''
کابو همدومره مهمه ښکاري۔ خو ``فرعي سټ دے'' د ``عنصر دے'' له مخې
\emph{تعريفولے} شو۔ د دې دپاره بايد د سټ تيورۍ په ژبه کښې داسې
!!a{formula}~$!A(x, y)$ ومومو چې د سټونو يوه جوړه~$\tuple{X, Y}$
هله او يوازې هله پوره کړي چې~$X \subseteq Y$ وي۔ خو~$X$ هله او
يوازې هله د~$Y$ فرعي سټ دے چې د~$X$ ټول !!{element}s د~$Y$ هم
!!{element}s وي۔ نو~$\subseteq$ په لاندې فارمول تعريفولے شو:
\[
\lforall[z][(z \in x \lif z \in y)]
\]
اوس چې هر کله په يوه فارمول کښې اړيکه~$\subseteq$ کارول غواړو، د
هغې پر ځاے همدا فارمول کارولے شو (په دې شرط چې~$x$ او~$y$ په مناسب
ډول بدل او تړلے متغير~$z$ د اړتيا پر وخت بيا ونومول شي)۔ مثلاً د
سټونو امتداديتوب وايي چې که هر دوه سټونه~$x$ او~$y$ په دوه اړخيزه
توګه په يو بل کښې شامل وي، نو~$x$ او~$y$ بايد هماغه يو سټ وي۔ دا په
$\lforall[x][\lforall[y][((x \subseteq y \land y
    \subseteq x) \lif x = y)]]$ بيانېدے شي؛ يا که~$\subseteq$ په
پورته تعريف بدل کړو، په لاندې ډول:
\[
\lforall[x][\lforall[y][((\lforall[z][(z \in x \lif z \in y)] \land
    \lforall[z][(z \in y \lif z \in x)]) \lif x = y)]].
\]
دا په رښتيا د~$\Log{ZFC}$ يو بديهي اصل، يعنې ``د امتداديتوب بديهي
اصل'' دے۔

د~$\emptyset$ دپاره هېڅ !!{constant} نشته، خو ``$x$ تش دے'' په
$\lnot \lexists[y][y \in x]$ بيانولے شو۔ بيا ``$\emptyset$ شته''
په !!{sentence}~$\lexists[x][\lnot \lexists[y][y \in x]]$ اوړي۔
دا د~$\Log{ZFC}$ يو بل بديهي اصل دے۔ (پام وکړئ چې د امتداديتوب
بديهي اصل دا لازموي چې يوازې يو تش سټ وي۔) هر کله چې د سټ تيورۍ په
ژبه کښې د~$\emptyset$ په باب خبرې کول غواړو، داسې به ليکو: ``داسې
يو سټ شته چې تش دے او~\dots''۔ د بېلګې په توګه، دا حقيقت چې
$\emptyset$ د هر سټ فرعي سټ دے، په لاندې ډول بيانولے شو:
\[
\lexists[x][(\lnot\lexists[y][y \in x] \land \lforall[z][x \subseteq
    z])]
\]
چېرته چې، البته، $x \subseteq z$ به بيا په خپل تعريف بدلېږي۔

د سټونو د عملونو، لکه~$X \cup Y$ او~$\Pow{X}$، په باب د خبرو دپاره
ورته چل کارول پکار دي۔ د سټ تيورۍ په ژبه کښې د تابعې سمبولونه نشته،
خو تابعوي اړيکې~$X \cup Y = Z$ او~$\Pow{X} = Y$ په لاندې ډول
بيانولے شو:
\begin{align*}
& \lforall[u][((u \in x \lor u \in y) \liff u \in z)]\\
& \lforall[u][(u \subseteq x \liff u \in y )]
\end{align*}
ځکه د~$X \cup Y$ !!{element}s کټ مټ هغه سټونه دي چې يا د~$X$
!!{element}s وي يا د~$Y$ (يا د دواړو) !!{element}s وي، او د~$\Pow{X}$
!!{element}s کټ مټ
د~$X$ فرعي سټونه دي۔ خو دا اجازه نۀ راکوي چې~$x \cup y$ يا~$\Pow{x}$
د ترمونو په توګه وکاروو: يوازې ټول هغه !!{formula}s کارولے شو چې
اړيکې~$X \cup Y = Z$ او~$\Pow{X} = Y$ تعريفوي۔ په اصل کښې لا دا هم
نۀ پوهېږو چې دغه اړيکې کله هم پوره کېږي، يعنې نۀ پوهېږو چې اتحادونه
او د فرعي سټونو سټونه تل شته۔ مثلاً !!{sentence}
$\lforall[x][\lexists[y][\Pow{x} = y]]$ د~$\Log{ZFC}$ يو بل بديهي
اصل دے (د فرعي سټونو د سټ بديهي اصل)۔

اوس د مرتبو جوړو يا تابعو په باب خبرې څنګه کوو؟ دلته بايد څرګنده کړو
چې مرتبې جوړې او تابعې څنګه د سټونو ځانګړي ډولونه ګڼلے شو۔ د مرتبې
جوړې~$\tuple{x, y}$ د تعريف يوه طريقه سټ~$\{\{x\}, \{x, y\}\}$ ده۔
خو لکه مخکښې، داسې !!a{function} نۀ شو معرفي کولے چې دغه سټ ونوموي؛
يوازې اړيکه~$\tuple{x, y} = z$، يعنې
$\{\{x\}, \{x, y\}\} = z$، تعريفولے شو:
\[
\lforall[u][(u \in z \liff (\lforall[v][(v \in u \liff v = x)] \lor
  \lforall[v][(v \in u \liff (v = x \lor v = y))]))]
\]
دا وايي چې د~$z$ !!{element}s~$u$ کټ مټ هغه سټونه دي چې يا يوازې
$x$ ئې !!{element} وي، يا يوازې~$x$ او~$y$ ئې !!{element}s وي (په
بل عبارت، هغه سټونه چې يا له~$\{x\}$ سره يو شان وي يا
له~$\{x, y\}$ سره)۔ چې دا ولرو، نورې خبرې هم کولے شو؛ مثلاً
$X \times Y = Z$:
\[
\lforall[z][(z \in Z \liff \lexists[x][\lexists[y][(x \in
        X \land y \in Y \land \tuple{x, y} = z)]])]
\]

تابعه~$f \colon X \to Y$ د اړيکې~$f(x) = y$ په توګه، يعنې د جوړو
د سټ~$\Setabs{\tuple{x,y}}{f(x) = y}$ په توګه، ګڼلے شو۔ بيا ويلے شو
چې سټ~$f$ له~$X$ څخه~$Y$ ته تابعه ده، که (الف) يوه اړيکه
$\subseteq X \times Y$ وي؛ (ب) ټوله وي، يعنې د هر~$x \in X$ دپاره
داسې~$y \in Y$ وي چې~$\tuple{x, y} \in f$؛ او (ج) تابعوي وي، يعنې
هر کله چې~$\tuple{x, y}, \tuple{x, y'} \in f$ وي، نو~$y = y'$ وي
(ځکه د تابعې قيمتونه بايد يوازيني وي)۔ نو ``$f$ له~$X$ څخه~$Y$ ته
تابعه ده'' داسې ليکل کېدے شي:
\begin{align*}
 \lforall[u][(u \in f \lif {} & \lexists[x][\lexists[y][(x \in X \land y \in
      Y \land \tuple{x, y} = u)]])] \land {}\\
 \lforall[x][(x \in X \lif {} &
      (\lexists[y][(y \in Y \land \mathrm{maps}(f, x, y))] \land {}\\
& (\lforall[y][\lforall[y'][((\mathrm{maps}(f, x, y) \land
    \mathrm{maps}(f, x, y')) \lif y = y')]])))]
\end{align*}
چېرته چې~$\mathrm{maps}(f, x, y)$ د
$\lexists[v][(v \in f \land \tuple{x, y} = v)]$ لنډون دے (دا
!!{formula} ``$f(x) = y$'' بيانوي)۔
\textbf{د ژباړې د نسخې سمون (OLFOL-026):} سرچينې په دواړو کلي
شرطونو کښې ساحوي قوسونه د شرط له پايلې وړاندې تړلي وو؛ دلته دواړه
پايلې د خپلو کميت ټاکونکو په ساحه کښې راوستل شوې دي۔

اوس دا هم ستونزمنه نۀ ده چې مثلاً بيان کړو چې~$f\colon X \to Y$
!!{injective} ده:
\begin{multline*}
 f \colon X \to Y \land \lforall[x][\lforall[x'][((x \in X \land x' \in
     X \land {} \\
 \lexists[y][(\mathrm{maps}(f, x, y) \land \mathrm{maps}(f,
        x', y))]) \lif x = x')]]
\end{multline*}
تابعه~$f\colon X \to Y$ هله او يوازې هله !!{injective} ده چې هر
کله~$f$، $x, x' \in X$ يوه~$y$ ته ولېږي، نو~$x = x'$ وي۔ که دا
فارمول په~$\mathrm{inj}(f, X, Y)$ لنډ کړو، لا له مخکښې په داسې مقام
کښې يو چې د سټ تيورۍ په ژبه د کانتور د قضيې غوندې نۀ معمولي خبره
بيان کړو: له~$\Pow{X}$ څخه~$X$ ته هېڅ !!{injective} تابعه نشته:
\[
\lforall[X][\lforall[Y][(\Pow{X} = Y \lif
    \lnot\lexists[f][\mathrm{inj}(f, Y, X)])]]
\]
\textbf{د ژباړې د نسخې سمون (OLFOL-027):} سرچينې دواړه کلي سوري د
مشترک قيمت د شتون له شرط وړاندې تړلي وو؛ دلته د هماغو سورو تړونکي
قوسونه د بشپړ شرطي فارمول تر پای وروسته راغلي دي۔

کېدے شي څوک فکر وکړي چې سټ تيوري يو بل داسې بديهي اصل غواړي چې د هر
تعريفوونکي خاصيت دپاره د يوه سټ شتون تضمين کړي۔ که~$!A(x)$ د سټ
تيورۍ داسې فارمول وي چې متغير~$x$ پکښې ازاد وي، لاندې !!{sentence}
په پام کښې نيولے شو:
\[
\lexists[y][\lforall[x][(x \in y \liff !A(x))]].
\]
دا !!{sentence} وايي چې داسې سټ~$y$ شته چې !!{element}s ئې ټول او
يوازې هغه~$x$ دي چې~$!A(x)$ پوره کوي۔ دا شېما ``د ادراک اصل'' بلل
کېږي۔ ډېره ګټوره ښکاري؛ خو له بده مرغه ناسازګاره ده۔
$!A(x) \ident \lnot x \in x$ واخلئ؛ بيا د ادراک اصل وايي:
\[
\lexists[y][\lforall[x][(x \in y \liff x \notin x)]],
\]
يعنې د ټولو هغو سټونو د سټ شتون وايي چې د خپل ځان !!{element}s نۀ
وي۔ داسې سټ نۀ شي موجودېدے---دا د راسل تناقض دے۔ په اصل کښې
$\Log{ZFC}$ د دې اصل يوه محدوده---او سازګاره---بڼه لري، يعنې د
بېلتون اصل:
\[
\lforall[z][\lexists[y][\lforall[x][(x \in y \liff (x \in z \land
      !A(x))]]].
\]

\begin{prob}
داسې !!a{derivation} ورکړئ چې لاندې وښيي، او په دې توګه ثابته کړئ
چې د ادراک اصل ناسازګار دے:
\[
\lexists[y][\lforall[x][(x \in y \liff x \notin x)]] \Proves \lfalse.
\]
ښايي لومړے د لاندې ثبوت ګټور وي:
$(A \lif \lnot A) \land (\lnot A \lif A) \Proves \lfalse$۔
\end{prob}
\end{document}
